\begin{proof}[Proof of \cref{obj:thm:sharp-interval-frontier}]
\begin{enumerate}
\item Since \(n\ge2\), the sampling term is strictly positive. Hence
\[
0<n^{-1/4}
\le
\max\left\{
n^{-1/4},(n^2\epsilon)^{-1/7},(n\epsilon)^{-1/2}
\right\},
\qquad
r(n,\epsilon)>0.
\]
The definitions of the benchmark and constants therefore give
\[
\ell(n,\epsilon)=r(n,\epsilon)>0,
\qquad
c_I=2^{-20}>0,
\qquad
C_I=2^{22}>0.
\]
% lean: CausalSmith.Stat.PrivateCateRoughdesign.rate_pos

\item We first identify the inversion kernel with the publicly tuned interval kernel. If \(r(n,\epsilon)\ge1/8\), both constructions in
\cref{obj:def:interval-handle,obj:def:public-tuning} return \([-1,1]\).

Suppose instead that \(r(n,\epsilon)<1/8\). Define the integer rounding index and public parameters by
\[
\jmath_h
=
\left\lceil\log_2\!\left(\frac1{r(n,\epsilon)}\right)\right\rceil
\in\mathbb Z_{\ge0},
\qquad
h=2^{-\jmath_h},
\qquad
k=2\,2^{4\jmath_h}\in\mathbb N.
\]
The ceiling inequalities imply
\[
\log_2\!\left(\frac1{r(n,\epsilon)}\right)
\le\jmath_h
\le
\log_2\!\left(\frac1{r(n,\epsilon)}\right)+1,
\]
and thus
\[
\frac1{r(n,\epsilon)}
\le2^{\jmath_h}
\le\frac2{r(n,\epsilon)},
\qquad
0<\frac{r(n,\epsilon)}2\le h\le r(n,\epsilon)<\frac18.
\]
In particular,
\[
0<h\le\frac14,
\qquad
k=2h^{-4}\ge2,
\qquad
\delta=\frac{2h}{k}=h^5>0.
\]
This also proves that the integer conversion in the public cell count is exact.

For every possible value of the joint release in
\cref{obj:def:private-ratio}, define
\[
\mathcal A=nh\min\{1,n\delta\},
\qquad
d_0=\frac{3\mathcal A}{64},
\qquad
B_D=\max\{d_0,\widetilde S_D\},
\qquad
B_N=\operatorname{clip}_{[-B_D,B_D]}(\widetilde S_N),
\qquad
\rho=\sqrt{10V(h,\delta)}.
\]
Because \(n,h,\delta>0\),
\[
\mathcal A>0,
\qquad
B_D\ge d_0>0,
\qquad
\rho\ge0.
\]
Division by the positive \(B_D\) preserves minima and maxima, so
\[
\frac{B_N}{B_D}
=
\frac{\max\{-B_D,\min\{B_D,\widetilde S_N\}\}}{B_D}
=
\max\left\{-1,\min\left\{1,\frac{\widetilde S_N}{B_D}\right\}\right\}
=
\widehat T_{h,k}.
\]
Consequently, the accepted candidate set is
\[
\begin{aligned}
\mathcal V_{\mathrm{acc}}
&=
\left\{\vartheta\in[-1,1]:
|B_N-\vartheta B_D|\le\rho B_D\right\}\\
&=
\left\{\vartheta\in[-1,1]:
\left|\frac{B_N}{B_D}-\vartheta\right|\le\rho\right\}\\
&=
\left[
\max\{-1,\widehat T_{h,k}-\rho\},
\min\{1,\widehat T_{h,k}+\rho\}
\right].
\end{aligned}
\]
Indeed, the absolute-value inequality is equivalent to
\(\widehat T_{h,k}-\rho\le\vartheta\le\widehat T_{h,k}+\rho\).
Moreover,
\[
-1\le\widehat T_{h,k}\le1,
\qquad
\max\{-1,\widehat T_{h,k}-\rho\}
\le\widehat T_{h,k}
\le
\min\{1,\widehat T_{h,k}+\rho\}.
\]
Thus the accepted set is a nonempty closed interval, and its hull has exactly these endpoints. The two constructions apply equal measurable interval maps to the same joint release. Together with the constant branch, this proves
\[
I^*_{n,\epsilon}=\widehat I^*.
\]
% lean: CausalSmith.Stat.PrivateCateRoughdesign.optimalIntervalHandle_eq_publicTunedInterval

\item By \cref{obj:lem:frontier-testing-priors}, choose positive integers \(s_0,s_1\), families
\[
(P^{(0)}_\iota)_{\iota=1}^{s_0}\subset\mathcal P,
\qquad
(P^{(1)}_\iota)_{\iota=1}^{s_1}\subset\mathcal P,
\]
and a separation satisfying
\[
\theta(P^{(0)}_\iota)=0,
\qquad
\theta(P^{(1)}_\iota)=\Delta,
\qquad
\Delta\ge2^{-14}r(n,\epsilon)>0.
\]
Their dataset mixture laws are
\[
Q^{(0)}
=
\frac1{s_0}\sum_{\iota=1}^{s_0}
\bigl((P^{(0)}_\iota)^{\mathrm{obs}}\bigr)^{\otimes n},
\qquad
Q^{(1)}
=
\frac1{s_1}\sum_{\iota=1}^{s_1}
\bigl((P^{(1)}_\iota)^{\mathrm{obs}}\bigr)^{\otimes n}.
\]
The testing-prior construction uses the bounded-differences log-mgf conclusion of
\citep[Theorem 6.2, proof, printed p.~167]{BoucheronLugosiMassart2013}
and the labelled-tree count of
\citep[Section 5.6, Theorem 5.40]{KellerTrotterAppliedCombinatorics}.
Its conclusion gives, for every private Markov interval kernel \(I\),
\[
d_{\mathrm{TV}}(IQ^{(0)},IQ^{(1)})\le\frac14.
\]

Fix any kernel \(I\) feasible for \cref{obj:def:interval-criterion}. Define its two output probability laws and the measurable containment events by
\[
\mathsf J_0=IQ^{(0)},
\qquad
\mathsf J_1=IQ^{(1)},
\qquad
\mathsf E_0=\{\mathsf i\in\mathcal I:0\in\mathsf i\},
\qquad
\mathsf E_1=\{\mathsf i\in\mathcal I:\Delta\in\mathsf i\}.
\]
The finite families are nonempty, so honesty at every constituent law yields
\[
\begin{aligned}
\mathsf J_0(\mathsf E_0)
&=
\frac1{s_0}\sum_{\iota=1}^{s_0}
\Pr_{((P^{(0)}_\iota)^{\mathrm{obs}})^{\otimes n},I}
\{0\in I(D)\}
\ge\frac9{10},\\
\mathsf J_1(\mathsf E_1)
&=
\frac1{s_1}\sum_{\iota=1}^{s_1}
\Pr_{((P^{(1)}_\iota)^{\mathrm{obs}})^{\otimes n},I}
\{\Delta\in I(D)\}
\ge\frac9{10}.
\end{aligned}
\]
Total variation comparison on \(\mathsf E_1\) gives
\[
\mathsf J_1(\mathsf E_1)-\mathsf J_0(\mathsf E_1)
\le d_{\mathrm{TV}}(\mathsf J_0,\mathsf J_1)
\le\frac14.
\]
Using the union--intersection identity and
\(\mathsf J_0(\mathsf E_0\cup\mathsf E_1)\le1\), we obtain
\[
\begin{aligned}
\mathsf J_0(\mathsf E_0\cap\mathsf E_1)
&=
\mathsf J_0(\mathsf E_0)+\mathsf J_0(\mathsf E_1)
-\mathsf J_0(\mathsf E_0\cup\mathsf E_1)\\
&\ge\frac9{10}+\left(\frac9{10}-\frac14\right)-1
=\frac{11}{20}.
\end{aligned}
\]

For every interval \(\mathsf i\in\mathsf E_0\cap\mathsf E_1\), its endpoint inequalities imply
\[
(0,\Delta)\subseteq\mathsf i,
\qquad
\operatorname{Leb}(\mathsf i)
\ge\operatorname{Leb}((0,\Delta))=\Delta.
\]
This argument applies to open and closed endpoints and to unbounded intervals; an empty interval cannot belong to the intersection. Therefore, with all length integrals understood as nonnegative extended integrals,
\[
\begin{aligned}
\frac{11}{20}\Delta
&\le\Delta\,\mathsf J_0(\mathsf E_0\cap\mathsf E_1)\\
&=
\int_{\mathcal I}
\Delta\,\mathbf 1_{\mathsf E_0\cap\mathsf E_1}(\mathsf i)
\,\mathsf J_0(d\mathsf i)\\
&\le
\int_{\mathcal I}\operatorname{Leb}(\mathsf i)\,
\mathsf J_0(d\mathsf i)\\
&=
\frac1{s_0}\sum_{\iota=1}^{s_0}
E_{((P^{(0)}_\iota)^{\mathrm{obs}})^{\otimes n},I}
\operatorname{Leb}(I(D))\\
&\le
\sup_{P\in\mathcal P}
E_{(P^{\mathrm{obs}})^{\otimes n},I}
\operatorname{Leb}(I(D)).
\end{aligned}
\]
The last inequality follows by bounding each constituent expectation by the same supremum.

The numerical calibration is
\[
2^{-20}\le\frac{11}{20}\,2^{-14}.
\]
Multiplying by the nonnegative benchmark and using the separation bound gives
\[
2^{-20}r(n,\epsilon)
\le
\frac{11}{20}\,2^{-14}r(n,\epsilon)
\le
\frac{11}{20}\Delta
\le
\sup_{P\in\mathcal P}
E_{(P^{\mathrm{obs}})^{\otimes n},I}
\operatorname{Leb}(I(D)).
\]
Since this holds for every feasible \(I\), taking the defining infimum proves
\[
2^{-20}r(n,\epsilon)\le H_{n,\epsilon}.
\]
% lean: CausalSmith.Stat.PrivateCateRoughdesign.finite_priors_interval_lower

\item Apply \cref{obj:thm:uniform-private-upper} with auxiliary cell count two and localization radius \(1/4\), which satisfy its parameter conditions. Its variance argument uses the replacement-copy Efron--Stein inequality from
\citep[Theorem 3.1, printed p.~54]{BoucheronLugosiMassart2013}.
The publicly tuned conclusions establish that \(\widehat I^*\) is an everywhere-defined Markov kernel satisfying pure replacement \(\epsilon\)-privacy, and that
\[
\Pr_{(P^{\mathrm{obs}})^{\otimes n},\widehat I^*}
\{\theta(P)\in\widehat I^*(D)\}
\ge\frac9{10}
\qquad(P\in\mathcal P),
\]
\[
\sup_{P\in\mathcal P}
E_{(P^{\mathrm{obs}})^{\otimes n},\widehat I^*}
\operatorname{Leb}(\widehat I^*(D))
\le2^{22}r(n,\epsilon).
\]
Thus \(\widehat I^*\) is feasible for the infimum in
\cref{obj:def:interval-criterion}, and
\[
H_{n,\epsilon}
\le
\sup_{P\in\mathcal P}
E_{(P^{\mathrm{obs}})^{\otimes n},\widehat I^*}
\operatorname{Leb}(\widehat I^*(D)).
\]
The kernel identity proved above transfers privacy, coverage, and length to \(I^*_{n,\epsilon}\), yielding
\[
2^{-20}r(n,\epsilon)
\le H_{n,\epsilon}
\le
\sup_{P\in\mathcal P}
E_{(P^{\mathrm{obs}})^{\otimes n},I^*_{n,\epsilon}}
\operatorname{Leb}(I^*_{n,\epsilon}(D))
\le2^{22}r(n,\epsilon).
\]
% lean: CausalSmith.Stat.PrivateCateRoughdesign.uniform_private_upper

\item Finally, \cref{obj:thm:matched-risk-frontier} supplies
\[
2^{-20}r(n,\epsilon)
\le R_{n,\epsilon}
\le2^{18}r(n,\epsilon).
\]
Multiplication by nonnegative finite constants preserves order in the extended nonnegative reals. The risk upper bound and interval lower bound give
\[
2^{-38}R_{n,\epsilon}
\le
2^{-38}\,2^{18}r(n,\epsilon)
=
2^{-20}r(n,\epsilon)
\le H_{n,\epsilon}.
\]
The interval upper bound and risk lower bound give
\[
H_{n,\epsilon}
\le2^{22}r(n,\epsilon)
=
2^{42}\bigl(2^{-20}r(n,\epsilon)\bigr)
\le2^{42}R_{n,\epsilon}.
\]
These are the asserted comparisons of the two frontiers.
% lean: CausalSmith.Stat.PrivateCateRoughdesign.sharp_interval_frontier
\end{enumerate}
\end{proof}
